\begin{pastexam}[label={px:nchu-114-dsa-02}]{中興資工所甲組 114 學年度 單選題第 2 題}
（單選；每題 3 分，未答題得 0 分，答錯不倒扣）

Which of the following sorting algorithms has the best worst-case time complexity?

(A)~Bubble Sort \qquad (B)~Merge Sort \qquad (C)~Insertion Sort \qquad (D)~Selection Sort
\end{pastexam}

\begin{pastsoln}
答案是 (B)。

四種排序的最壞情況（資結書第~14.2、14.4、14.7~節）：
\begin{itemize}
\item (A) Bubble Sort：輸入由大到小時每一趟都要交換，比較 $n(n-1)/2$ 次，$\Theta(n^2)$。
\item (B) Merge Sort：$T(n) = 2T(n/2) + \Theta(n)$，不論輸入長什麼樣子都是 $\Theta(n \log n)$。
\item (C) Insertion Sort：輸入由大到小時第 $i$ 個要往前移 $i - 1$ 格，共 $n(n-1)/2$ 次，$\Theta(n^2)$。
\item (D) Selection Sort：每一趟都要掃完剩下的部分找最小值，任何輸入都比較 $n(n-1)/2$ 次，$\Theta(n^2)$。
\end{itemize}
只有 Merge Sort 的最壞情況是 $\Theta(n \log n)$，而這也是比較式排序最壞情況的下界，選 (B)。
\end{pastsoln}
